\renewcommand{\thesection}{S\arabic{section}}
\setcounter{section}{4}
\section{Controlled synthetic event-generation framework}
\label{sec:Appendix_S5}

This section describes the controlled synthetic framework used to construct the event-identifiability landscapes in Fig.~\ref{fig:event-impact}. The objective of these simulations was not to reproduce a specific historical epidemic, but to isolate how different event parameterizations propagate through epidemic dynamics, delay filtering, and observational noise under a common epidemic setting.

\subsection{Baseline epidemic}

All event classes were evaluated using the same baseline epidemic trajectory. The baseline incidence curve was defined as a smooth bimodal epidemic,

\begin{equation}
U_0(t)
=
35
+
105
\exp\!\left[
-\frac{(t-42)^2}{2\times16^2}
\right]
+
85
\exp\!\left[
-\frac{(t-88)^2}{2\times23^2}
\right],
\end{equation}

evaluated over a simulation window consisting of 30 pre-event days and 120 post-event days. The event was centered at day 30.

The corresponding baseline reproduction-number trajectory was inferred from the renewal equation,

\begin{equation}
R_0(t)
=
\frac{U_0(t)}
{\sum_k w_kU_0(t-k)},
\end{equation}

where \(w_k\) denotes the discrete generation-interval distribution. This baseline \(R_t\) trajectory served only as the reference transmission process on which event perturbations were imposed.



\subsection{Generation interval}

Transmission dynamics followed the discrete renewal model described in the main text. The generation interval was represented by a Gamma distribution with mean 5 days and standard deviation 2 days, discretized over lags of 1--21 days and normalized to sum to one.



\subsection{Event parameterization}

Each event modified the baseline reproduction number through

\begin{equation}
R(t;\theta)
=
R_0(t)m(t;\theta),
\end{equation}

where

\[
\theta=(a,\tau)
\]

contains the peak proportional change in \(R_t\) (\(a\)) and the event duration (or rollout timescale) (\(\tau\)).

Three stylized transmission perturbations were considered.

\paragraph{Mass gathering.}

A transient increase in transmission was represented as

\[
m(t;\theta)
=
1
+
a
h_\tau(t-t_0),
\]

where \(h_\tau\) is a smoothed finite-duration window normalized to have unit maximum.

\paragraph{Temporary lockdown.}

A transient reduction in transmission was represented as

\[
m(t;\theta)
=
1
-
a
h_\tau(t-t_0),
\]

using the same normalized event window.

\paragraph{Vaccination rollout.}

Vaccination was represented by a monotone rollout curve,

\[
m(t;\theta)
=
1
-
a
C_\tau(t),
\]

where \(C_\tau(t)\) is a normalized logistic coverage curve increasing from 0 to 1 over the rollout period.

The event window \(h_\tau\) was constructed using two opposing logistic transitions,

\[
h_\tau(t)
=
\sigma\!\left(
\frac{t-t_{\rm start}}{e}
\right)
-
\sigma\!\left(
\frac{t-t_{\rm end}}{e}
\right),
\]

where

\[
\sigma(x)
=
\frac1{1+e^{-x}},
\]

followed by normalization to unit maximum.

The vaccination rollout curve was generated using a logistic function centered at the midpoint of the rollout period and linearly rescaled so that

\[
C_\tau(t_{\rm start})=0,
\qquad
C_\tau(t_{\rm end})=1.
\]



\subsection{Renewal simulation}

For every candidate event profile,

\[
R(t;\theta)
=
R_0(t)m(t;\theta),
\]

the corresponding incidence trajectory was generated recursively using

\begin{equation}
U(t;\theta)
=
R(t;\theta)
\sum_k
w_k
U(t-k;\theta).
\end{equation}

The first generation-interval window was initialized using the baseline trajectory. Consequently, within each event family every candidate profile was generated by perturbing the same baseline epidemic through the same renewal process. Differences across the event-identifiability landscape therefore arose solely from differences in the event parameterization under a common epidemic and observation setting.



\subsection{Observation model}

The downstream conditional mean was obtained through convolution with a common reporting-delay distribution,

\[
D(t;\theta)
=
p
(U*g)(t),
\]

where \(p=1\) in all controlled simulations.

The delay distribution was modeled by a discretized lognormal distribution with mean 7 days, standard deviation 6 days, and maximum delay 40 days.

Full linear convolution was retained throughout,

\[
D
=
U*g,
\]

rather than truncating the convolution to the original observation window. This preserves the exact Fourier relationship

\[
\widehat D
=
\widehat G
\widehat U
\]

after zero-padding and ensures exact agreement between the time-domain and spectral implementations of the separability functional.



\subsection{Negative-binomial observations}

Illustrative observations were generated independently from

\[
Y_t
\sim
\mathrm{NB}
(\mu_t,\kappa),
\]

where

\[
\mu_t=D(t;\theta),
\]

and

\[
\mathrm{Var}(Y_t)
=
\mu_t
+
\frac{\mu_t^2}{\kappa}.
\]

A common dispersion parameter

\[
\kappa=20
\]

was used throughout all simulations.

One synthetic realization was generated from each reference event and displayed in Fig.~\ref{fig:event-impact} only to illustrate the scale of observational variability. All identifiability calculations were performed using the corresponding downstream conditional means rather than the sampled observations.



\subsection{Reference event profiles}

The reference event profiles corresponded to a peak proportional \(50\%\) change in \(R_t\).

The reference durations were

\[
3~\text{days},
\qquad
10~\text{days},
\qquad
60~\text{days},
\]

for the mass gathering, temporary lockdown, and vaccination rollout, respectively.

Candidate event landscapes were evaluated over

\[
a\in[0.05,0.90]
\]

and event-specific duration grids covering

\[
1\!-\!30,
\qquad
3\!-\!60,
\qquad
10\!-\!120
\]

days for the three event classes.



\subsection{Construction of the event-identifiability landscapes}

For every point on the parameter grid,

\[
\theta=(a,\tau),
\]

the corresponding incidence trajectory was generated through the renewal model and mapped to the downstream conditional mean.

The event-identifiability landscape was then computed using

\[
J_{\rm event}
(\theta,\theta^\ast)
=
\int
\frac{
|p\widehat G(f)|^2
|\Delta\widehat U(f)|^2
}
{N(f)}
\,df,
\]

where

\[
\Delta U
=
U(\theta)-U(\theta^\ast).
\]

Thus every point in the landscape represents the downstream separability between a candidate event profile and the specified reference profile.



\subsection{Selection of illustrative alternatives}

For each event class, one representative alternative was selected automatically for visualization.

Candidates were restricted to those lying immediately outside the sufficient non-identifiability region,

\[
J_{\rm threshold}
<
J_{\rm event}
\le
3J_{\rm threshold},
\]

while satisfying minimum parameter separation from the reference and excluding points near the parameter-space boundary.

Among the remaining candidates, the displayed example was chosen to maximize the relative upstream root-mean-square difference from the reference trajectory. This procedure produces examples with substantial differences in transmission dynamics while remaining close to the theoretical non-identifiability boundary.


